\section{总结与延伸}

\subsection{核心结论}

\begin{enumerate}
\item 妙鸭是把 AI 能力用得很好的互联网产品，但不是张月光后来定义的 AI Native。
\item AI Native 的关键变化，是从流程设计转向上下文设计，产品经理要组织模型可见的信息、工具和反馈。
\item One Way Door 是判断长期正确产品的重要信号：用户一旦迁移，就不愿回到旧解法。
\item AI 产品组织从线性工作变成两段式：先验证模型能力和 context，再封装流程和体验。
\item 张月光把这代 AI 理解为创造 AI 人口，而不是只创造 AI 服务。
\item Will / Skill 分工说明，人负责目标、判断和品味，AI 负责执行和能力扩展。
\item 星眠探索 AI 内容和陪伴，Dokie 探索创作表达 Agent 和能力边界突破。
\item Agent 产品可以先走短程高频反馈路线，而不是只追逐超长程离线任务。
\item LM 仍是当前 Agent 的核心内核，但这与物理世界智能是否需要超越语言是两个不同问题。
\end{enumerate}

\subsection{开放问题}

最后保留开放问题，是因为本期给出的更多是产品方法论和创业判断，而不是已经被市场完全验证的结论。下面这些问题，决定 Dokie 和同类 Agent 产品能否从工具走向入口。

\begin{itemize}
\item AI Native 产品是否会形成稳定的方法论，还是每一类产品都要重新发明 context design？
\item Agent 的最佳切入点是创作表达、资料处理、代码、办公自动化，还是人格化 AI 朋友？
\item AI 乙女游戏能否成为大品类，还是只会是 AI 内容的一条窄路径？
\item LM 作为 Agent 内核能持续多久，多模态模型会不会改变这个判断？
\item 用户会为“AI 朋友”付费，还是只会为“能稳定干活的工具”付费？
\end{itemize}

\subsection{拓展阅读}

\begin{itemize}
\item EP139 Agent 技术史：理解 Agent 从工具调用到社会辐射的脉络。
\item EP135 Tristan / Elys：理解赛博分身、context 流动和 AI 社交网络。
\item EP113 杨植麟访谈：Agentic LLM、缸中之脑和无限开端。
\item Cursor、Claude Code 等 coding agent 案例：理解 One Way Door 产品。
\item AI 游戏、互动叙事、Live2D 和虚拟陪伴相关材料：理解星眠的内容路线。
\end{itemize}

\begin{importantbox}{最后的判断}
本期最值得带走的不是“妙鸭不是 AI Native”这句反直觉判断，而是它背后的方法论：当输入和输出都开放时，产品不再只是流程；它变成了上下文、模型、工具、反馈、品味和用户心智共同构成的系统。
\end{importantbox}

\end{document}
